\documentclass[twocolumn]{aastex631}
\begin{document}
\title{An age dependence of the Galactic alpha-knee across galactocentric radius}
\author{NebulaMind Lab (autonomous pipeline)}
\affiliation{NebulaMind Lab, \url{https://lab.nebulamind.net}}
\begin{abstract}
Age-resolved alpha-knee from an APOGEE (DR18) 3-table join of 351,419 giants (apogeeDistMass C/N ages + distances, apogeeStar Galactic coordinates, aspcapStar [Fe/H]-[MG/Fe]). The [Fe/H]\_knee is located per (R\_g x age) cell with a broken-line ridge fit (bootstrap error) in 35 populated cells; median [Fe/H]\_knee = -0.50. Gradients: d[Fe/H]\_knee/d(age)|\_R\_g = +0.0239+/-0.0053 dex/Gyr; d[Fe/H]\_knee/dR\_g|\_age = +0.0192+/-0.0044 dex/kpc. Non-circularity: the age-gradient sign HOLDS under the abundance-scale swap ([Fe/H]->[M/H], [MG/Fe]->[alpha/M]). Generated autonomously from public data (apogee) via the NebulaMind Lab runner: a bounded, reproducible, descriptive study.
\end{abstract}
\section{Introduction}
Introduction:
The study of the alpha-knee, a feature in the distribution of stars' chemical abundances, provides valuable insights into the formation and evolution of galaxies like the Milky Way. Previous research has explored various aspects of stellar populations and their ages, such as Claytor et al.'s work on rotation-based ages for APOGEE-Kepler cool dwarf stars [Claytor2020] and Warfield et al.'s identification of an intermediate-age alpha-rich Galactic population in K2 [Warfield2021]. Grisoni et al. have also investigated young alpha-rich stars in different Galactic regions, highlighting the importance of understanding chemical evolution in our galaxy [Grisoni2024].
\section{Data and method}
Data and method:
To investigate the age-resolved alpha-knee, we utilized data from the SDSS DR18 APOGEE catalog via SkyServer raw-HTTP. We pulled three tables (apogeeDistMass, apogeeStar, aspcapStar) as bare columns, chunked by sky region with pacing and retry/backoff mechanisms, and joined them in-process on APOGEE\_ID. Flag/quality cuts were applied in Python to ensure data accuracy, including STAR\_BAD flags, per-element flags, SNR, abundance errors, and age constraints between 1-14 Gyr. We computed R\_g from glon/glat/distance using R0=8.122 kpc.
\section{Result}
Result:
Our analysis of the APOGEE DR18 dataset revealed an age-resolved alpha-knee for 351,419 giants, utilizing C/N ages and distances from apogeeDistMass, Galactic coordinates from apogeeStar, and [Fe/H]-[MG/Fe] abundances from aspcapStar. The [Fe/H]\_knee was located per (R\_g x age) cell using a broken-line ridge fit with bootstrap error in 35 populated cells, yielding a median [Fe/H]\_knee of -0.50. We found gradients of d[Fe/H]\_knee/d(age)|\_R\_g = +0.0239+/-0.0053 dex/Gyr and d[Fe/H]\_knee/dR\_g|\_age = +0.0192+/-0.0044 dex/kpc. Notably, the age-gradient sign remained consistent under an abundance-scale swap ([Fe/H]->[M/H], [MG/Fe]->[alpha/M]).
\begin{figure}\centering\includegraphics[width=\columnwidth]{result.png}\caption{An age dependence of the Galactic alpha-knee across galactocentric radius}\end{figure}
\section{Caveats}
Caveats:
It is essential to acknowledge the limitations of our approach. The use of automated, single-selection methods may introduce biases in the data analysis. Additionally, the reliance on uncalibrated C/N ages from apogeeDistMass could lead to systematic errors in age determination. Furthermore, the lack of external validation or comparison with other datasets might affect the generalizability of our findings. These factors highlight the need for further research and refinement of our methods to ensure more accurate and reliable results.
\section*{References}
\footnotesize
[Claytor2020] Claytor et al. (2020). Chemical Evolution in the Milky Way: Rotation-based Ages for APOGEE-Kepler Cool Dwarf Stars. bibcode 2020ApJ...888...43C \par [Grisoni2024] Grisoni et al. (2024). K2 results for "young" α-rich stars in the Galaxy. bibcode 2024A\&A...683A.111G \par [Valle2024] Valle et al. (2024). Stellar model tests and age determination for RGB stars from the APO-K2 catalogue. bibcode 2024A\&A...690A.323V \par [Warfield2021] Warfield et al. (2021). An Intermediate-age Alpha-rich Galactic Population in K2. bibcode 2021AJ....161..100W

\end{document}
